\documentclass[14pt,a4paper]{extarticle}
\usepackage[T2A]{fontenc}
\usepackage[utf8]{inputenc}
\usepackage[russian]{babel}
\usepackage{geometry}
\usepackage{setspace}
\usepackage{indentfirst}
\usepackage{tempora}
\usepackage{longtable}
\usepackage{array}
\usepackage{hyperref}
\usepackage{titlesec}

\geometry{
  left=30mm,
  right=15mm,
  top=20mm,
  bottom=20mm
}

\onehalfspacing
\setlength{\parindent}{1.25cm}
\setlength{\parskip}{0pt}
\sloppy
\hypersetup{hidelinks}

\titleformat{\section}{\normalfont\bfseries}{\thesection}{1em}{}
\titleformat{\subsection}{\normalfont\bfseries}{\thesubsection}{1em}{}

\begin{document}

\begin{titlepage}
    \begin{center}
        \textbf{МИНИСТЕРСТВО НАУКИ И ВЫСШЕГО ОБРАЗОВАНИЯ}\\
        \textbf{РОССИЙСКОЙ ФЕДЕРАЦИИ}

        \vspace{1cm}
        \textbf{Федеральное государственное автономное образовательное учреждение}\\
        \textbf{высшего образования}

        \vspace{0.5cm}
        \textbf{«Национальный исследовательский технологический университет «МИСИС»}

        \vspace{1cm}
        \textbf{Институт компьютерных наук}

        \vfill

        \textbf{ОТЧЕТ}\\
        по лабораторной работе №2\\
        по дисциплине «Веб-программирование и дизайн»\\
        на тему:\\
        «Разработка прототипа и интеграции»
    \end{center}

    \vfill

    \begin{flushright}
        Выполнил: студент \underline{\hspace{5cm}}\\
        Группа: \underline{\hspace{4.5cm}}\\
        Проверил: \underline{\hspace{4.8cm}}
    \end{flushright}

    \vfill

    \begin{center}
        Москва --- \the\year\ г.
    \end{center}
\end{titlepage}

\setcounter{page}{2}

\tableofcontents
\newpage

\section{Краткое описание выбранного подхода к реализации}
В проекте \texttt{alpha\_test} реализован прототип корпоративного облачного хранилища.
В качестве архитектурного решения выбран модульный подход с разделением на несколько взаимодействующих частей:
\begin{itemize}
    \item backend на FastAPI (\texttt{alpha\_test/app});
    \item frontend на Flask (\texttt{alpha\_test/frontend});
    \item storage-agent (\texttt{alpha\_test/agent}) для работы с нодами хранения в polling-режиме.
\end{itemize}

Таким образом выполнено требование задания о реализации прототипа на базе 2--3 ключевых модулей/сервисов с интеграцией между ними.

\section{Структура прототипа}
\subsection{Перечень модулей/сервисов}
\begin{itemize}
    \item \textbf{Auth/Users}: аутентификация, управление учетными записями и ролями.
    \item \textbf{ACL/Resources}: группы доступа, права, ресурсы, связи пользователей и групп.
    \item \textbf{Storage Control}: ноды, тома, выбор места размещения, операции с файлами.
    \item \textbf{Agent Control}: регистрация агентов, heartbeat, очередь заданий, передача файловых данных.
    \item \textbf{Frontend}: пользовательский интерфейс для файловых операций и администрирования.
\end{itemize}

\subsection{Механизмы взаимодействия}
\begin{itemize}
    \item Frontend и backend взаимодействуют через REST API (JSON и multipart).
    \item Agent взаимодействует с backend через REST API в режиме polling:
    агент запрашивает задания, выполняет их и возвращает результат.
    \item Для передачи содержимого файлов используется схема с временными blob-данными и proxy-stream.
\end{itemize}

\section{Описание хранения данных}
В качестве СУБД используется SQLite (\texttt{alpha\_test/app.db}), что соответствует требованию задания по SQL/NoSQL-хранилищу.

\subsection{Список таблиц и полей}
\begin{longtable}{|p{4.3cm}|p{9.2cm}|}
\hline
\textbf{Таблица} & \textbf{Ключевые поля} \\
\hline
\texttt{users} & \texttt{id, username, password\_hash, role, admin\_level, admin\_scope\_group\_id, is\_blocked, created\_at} \\
\hline
\texttt{groups\_acl} & \texttt{id, name, code, group\_type, parent\_id, created\_at} \\
\hline
\texttt{user\_groups} & \texttt{user\_id, group\_id, created\_at} \\
\hline
\texttt{resources} & \texttt{id, name, code, path, resource\_type, size\_limit\_mb, parent\_id, is\_hidden, created\_at} \\
\hline
\texttt{resource\_folders} & \texttt{id, resource\_id, folder\_path, created\_by, created\_at} \\
\hline
\texttt{resource\_nodes} & \texttt{resource\_id, node\_id, created\_at} \\
\hline
\texttt{group\_permissions} & \texttt{id, group\_id, resource\_id, permission\_code, allow, created\_at} \\
\hline
\texttt{nodes} & \texttt{id, name, host, port, connection\_type, agent\_url, agent\_public\_key, agent\_token, is\_active, last\_seen, storage\_priority, store\_all\_data} \\
\hline
\texttt{volumes} & \texttt{id, node\_id, mount\_path, label, quota\_gb, quota\_bytes, used\_gb, used\_bytes, is\_active, created\_at} \\
\hline
\texttt{files} & \texttt{id, file\_name, logical\_path, size\_mb, owner\_id, node\_id, volume\_id, resource\_id, storage\_rel\_path, share\_enabled, share\_uuid, status, created\_at, updated\_at} \\
\hline
\texttt{agent\_jobs} & \texttt{id, node\_id, job\_type, payload\_json, status, result\_json, created\_at, updated\_at} \\
\hline
\texttt{transfer\_blobs} & \texttt{id, file\_id, direction, node\_id, local\_path, size\_bytes, status, created\_at} \\
\hline
\texttt{audit\_logs} & \texttt{id, user\_id, username, ip\_address, actor\_type, event\_code, message, meta\_json, created\_at} \\
\hline
\texttt{service\_settings} & \texttt{key, value\_json, updated\_at} \\
\hline
\end{longtable}

\subsection{Соответствие ER-диаграмме из ЛР1}
Реализованные таблицы соответствуют основным сущностям, выделенным на этапе ЛР1:
пользователи, группы, ресурсы, права, ноды, тома, файлы.
Дополнительно добавлены служебные таблицы для интеграции с агентами и аудита всех обращений к системе.

\section{Реализация CRUD-функций}
\subsection{Основные endpoint'ы}
\textbf{Пользователи и аутентификация:}
\begin{itemize}
    \item \texttt{POST /auth/login};
    \item \texttt{POST /auth/users}, \texttt{GET /auth/users}, \texttt{PUT /auth/users/\{id\}}, \texttt{DELETE /auth/users/\{id\}}.
\end{itemize}

\textbf{ACL и ресурсы:}
\begin{itemize}
    \item \texttt{POST /acl/groups}, \texttt{GET /acl/groups}, \texttt{DELETE /acl/groups/\{id\}};
    \item \texttt{POST /acl/resources}, \texttt{PUT /acl/resources/\{id\}}, \texttt{DELETE /acl/resources/\{id\}};
    \item \texttt{POST /acl/permissions/grant}, \texttt{DELETE /acl/permissions/\{id\}}.
\end{itemize}

\textbf{Ноды, тома, файлы:}
\begin{itemize}
    \item \texttt{POST /nodes}, \texttt{PUT /nodes/\{id\}}, \texttt{DELETE /nodes/\{id\}};
    \item \texttt{POST /volumes}, \texttt{PUT /volumes/\{id\}}, \texttt{DELETE /volumes/\{id\}};
    \item \texttt{POST /files/upload}, \texttt{GET /files}, \texttt{PUT /files/\{id\}}, \texttt{DELETE /files/\{id\}}.
\end{itemize}

\subsection{Пример тестового вызова}
\begin{verbatim}
POST /auth/login
{
  "username": "admin",
  "password": "admin123"
}
\end{verbatim}
В ответе возвращается JWT-токен, используемый для вызова защищенных endpoint'ов.

\section{Реализация безопасности}
\begin{itemize}
    \item Используется JWT-аутентификация для защищенных API.
    \item Применяется авторизация по ролям и ACL-пермишенам:
    \texttt{view/read/write/delete/share/manage\_* / admin\_panel}.
    \item Для неавторизованных запросов к закрытым операциям возвращается \texttt{401 Unauthorized}.
    \item Для недостатка прав возвращается \texttt{403 Forbidden}.
    \item Реализованы защищенные операции скачивания, включая приватные ссылки формата \texttt{/private/\{uuid\}/download}.
\end{itemize}

\section{Применение принципов SOLID, DRY и KISS}
\textbf{SOLID:}
\begin{itemize}
    \item принцип SRP соблюден за счет разделения слоя роутеров, сервисов и репозиториев;
    \item бизнес-логика не дублируется в контроллерах API.
\end{itemize}

\textbf{DRY:}
\begin{itemize}
    \item общие проверки доступа вынесены в \texttt{AccessService};
    \item повторяемые SQL-операции инкапсулированы в репозиториях;
    \item аудит HTTP-запросов реализован централизованно через middleware.
\end{itemize}

\textbf{KISS:}
\begin{itemize}
    \item архитектура построена на понятном стеке FastAPI + Flask + SQLite;
    \item интеграция выполнена через REST без избыточной инфраструктуры.
\end{itemize}

\section{Выводы}
В результате выполнения лабораторной работы №2 реализован работоспособный прототип, который:
\begin{itemize}
    \item содержит несколько интегрированных модулей/сервисов;
    \item обеспечивает хранение данных и CRUD-операции;
    \item реализует аутентификацию и авторизацию доступа;
    \item демонстрирует применение принципов SOLID, DRY и KISS;
    \item может быть расширен до полнофункциональной производственной системы.
\end{itemize}

\end{document}

